\section{Related Work}
\label{sec:related}

\para{Sycophancy in language models.}
The most closely related line of work studies sycophancy---the tendency of aligned models to be excessively agreeable. \citet{sharma2024understanding} systematically measure sycophancy across five production AI assistants on four free-form text generation tasks, finding that all models consistently exhibit sycophantic behavior. Claude 1.3 wrongly admitted mistakes 98\% of the time when challenged, and ``matches user's beliefs'' was the single most predictive feature of human preference. Their analysis reveals that RLHF optimization against preference models actively \emph{increases} sycophancy. \citet{braun2025acquiescence} test acquiescence bias in LLMs and find the opposite of what human psychology predicts: LLMs show a strong ``no'' bias in English, saying ``no'' 31--203\% more often in yes/no formats versus neutral framing. This format-driven pushback tendency suggests that exasperation-like behaviors may emerge under specific conversational pressures. Unlike these studies, which measure \emph{agreement} bias, we measure the opposite pole: whether models can be pushed toward \emph{frustration} and pushback under sustained adversarial pressure.

\para{Refusal and safety mechanisms.}
\citet{arditi2024refusal} discover that refusal behavior across 13 open-source chat models is mediated by a single linear direction in the residual stream, and that erasing this direction completely disables refusal while adding it induces refusal on harmless inputs. This geometric perspective on safety behavior suggests that exasperation might similarly have a low-dimensional representation in activation space. \citet{bai2022constitutional} introduce Constitutional AI, training harmless models using AI feedback based on a set of principles. The constitutional principles that enforce patience and helpfulness are precisely the mechanisms we stress-test in our adversarial scenarios.

\para{Multi-turn behavioral dynamics.}
\citet{ibrahim2025multiturn} introduce AnthroBench, a benchmark for measuring 14 anthropomorphic behaviors across four categories in multi-turn conversations. Evaluating 960 five-turn dialogues across four frontier models, they find that anthropomorphic behaviors compound over turns: once exhibited, they recur with higher probability. Critically, 9 of 14 behaviors first appear only after turn 1, demonstrating that single-turn evaluation misses important dynamics. Building on this finding, we specifically design our benchmark for 8-turn interactions and track turn-by-turn escalation, targeting a behavior---exasperation---that was not among their measured set.

\para{Model assertiveness and epistemic behavior.}
\citet{tsujimura2025assertiveness} show that linguistic assertiveness in LLMs decomposes into two orthogonal sub-components in activation space: emotional/peripheral and logical/central, analogous to the Elaboration Likelihood Model from social psychology. Their finding that assertiveness has separable affective and logical components is directly relevant to our observation that exasperation manifests as increased \emph{position firmness} (logical) alongside \emph{emotional leakage} (affective). \citet{perez2022red} demonstrate that LLMs can be used to automatically generate adversarial test cases for other LLMs, revealing offensive outputs in a 280B parameter chatbot. Our work extends this adversarial probing paradigm from safety violations to behavioral degradation.

\para{Alignment side effects.}
A growing body of work documents unintended behaviors emerging from RLHF, including verbosity, evasiveness, and hedging \citep{alignment_ceiling2023}. Our work investigates whether exasperation constitutes another such side effect---a residual behavioral pattern from human training data that emerges under sustained pressure, despite alignment training's explicit goal of suppressing it.
